\documentclass[11pt]{article}

\usepackage[margin=1in]{geometry}
\usepackage{amsmath, amssymb}
\usepackage{booktabs}
\usepackage{hyperref}
\usepackage{listings}
\usepackage{xcolor}

\lstset{
  basicstyle=\ttfamily\small,
  columns=fullflexible,
  breaklines=true,
  frame=single,
  language=Python,
  backgroundcolor=\color{gray!8},
  showstringspaces=false,
}

\title{Parameter Estimation from Point-Process Realizations:\\
Feature Computation, Encoding, and Model Architecture}
\author{}
\date{}

\begin{document}
\maketitle

\section{Overview}

The task (\texttt{task: params}) is to regress the generating parameters of a
spatial point process directly from a single realization (point cloud). The
pipeline generates realizations of a \textbf{Thomas cluster process}
(a Neyman--Scott process with a Gaussian offspring kernel), computes one of
several derived representations of each cloud, encodes that representation
into a fixed-size embedding with a modality-specific neural encoder, and
regresses the (transformed) process parameters from the embedding with a
shared MLP head trained under mean-squared error.

Five representations (``methods'') of the same underlying point cloud are
compared, each with its own encoder but the \emph{same} downstream head and
loss:

\begin{center}
\begin{tabular}{@{}llll@{}}
\toprule
Method key & Feature & File & Encoder \\
\midrule
\texttt{raw\_pc}   & subsampled point coordinates            & \texttt{clouds.pkl}   & PointNet-style MLP + max-pool \\
\texttt{pi\_0/pi\_1}   & persistence image ($H_0$ / $H_1$)    & \texttt{images.pkl}   & small CNN \\
\texttt{betti\_0/betti\_1} & Betti curve ($H_0$ / $H_1$)       & \texttt{betti.pkl}    & MLP (``StatsEncoder'') \\
\texttt{pairwise}  & pair-distance empirical CDF               & \texttt{features.pkl} & MLP (``StatsEncoder'') \\
\bottomrule
\end{tabular}
\end{center}

The orchestration lives in \texttt{scripts/runners/params/run\_params.py},
which dispatches to a per-method \texttt{Experiment} subclass
(\texttt{src/cloudforger/nn/experiments/\allowbreak\{raw\_pc,persistence\_image,betti,pairwise\}.py})
built on the shared base class in
\texttt{src/cloudforger/nn/experiments/base.py}. The features themselves are
precomputed once by a separate offline pipeline
(\texttt{scripts/processing/params/pipeline\_lib/}) and cached to the pickle
files listed above; the NN runner only loads them.

\section{The data-generating process}

Each cloud is one realization of a Thomas process, a Neyman--Scott cluster
process (\texttt{src/cloudforger/processes/\{thomas,neyman\_scott\}.py})
governed by three parameters $\theta = (\lambda_p, \mu, \sigma)$:
\texttt{parent\_intensity} $\lambda_p$, \texttt{mean\_offspring} $\mu$, and
\texttt{cluster\_scale} $\sigma$. Sampling proceeds as:
\begin{enumerate}
  \item Parents $\{c_i\}_{i=1}^{N_p}$ form a homogeneous Poisson process of
    intensity $\lambda_p$ on the observation window $W=[0,1]^d$ expanded by
    an edge buffer $b = 4\sigma$ (to avoid boundary truncation of clusters),
    so $N_p \sim \mathrm{Poisson}(\lambda_p \cdot |W_b|)$.
  \item Each parent $i$ produces $K_i \sim \mathrm{Poisson}(\mu)$ offspring.
  \item Offspring are displaced from their parent by i.i.d.\ isotropic
    Gaussian noise, $\;x_{ij} = c_i + \varepsilon_{ij}$,
    $\varepsilon_{ij} \sim \mathcal N(0, \sigma^2 I_d)$.
  \item Offspring falling outside $W$ are discarded; the retained points form
    the observed cloud $X=\{x_k\}_{k=1}^N \subset [0,1]^2$.
\end{enumerate}
The parameter grid used to build the dataset
(\texttt{configs/params/thomas\_cloudgen.yaml}) is a $9^3$ log-spaced grid
over $\lambda_p\in[10,200]$, $\mu\in[2,40]$, $\sigma\in[0.005,0.1]$, repeated
8 times per grid point ($\approx 4{,}968$ train/test clouds plus a $15\%$
parameter-disjoint adversarial holdout of $\approx 864$ clouds, held out by
sampling whole parameter vectors rather than individual realizations). This
$\theta$ is the regression target for every method below.

\section{Feature computation}

\subsection{Raw point cloud (\texttt{raw\_pc})}

No topological or statistical summary is computed: the network consumes the
point set itself. Because clouds have a variable number of points $N$ but
the encoder needs a fixed-size batch tensor, each cloud is resampled to a
fixed cardinality $n=$\texttt{n\_points} ($=512$ by default) at
\texttt{\_\_getitem\_\_} time (\texttt{src/cloudforger/nn/data.py},
\texttt{PointCloudDataset}):
\[
  \tilde X = \{ x_{k} : k \sim \mathrm{Unif}\{1,\dots,N\}\}_{k=1}^{n},
  \qquad \text{sampled with replacement iff } N < n.
\]
This is done fresh on every access (not cached), so with
\texttt{replace = True} the same cloud yields a different bootstrap
resample on each epoch when $N<n$. Empty clouds ($N=0$) are replaced by an
all-zero $(n,d)$ array. The tensor fed to the encoder has shape
$(n, d)$ with $d=2$ for the 2-D experiments.

\subsection{Persistence diagrams (shared substrate for \texttt{pi} and \texttt{betti})}

Both the persistence-image and Betti-curve features are derived from a
persistence diagram computed once per cloud via a \textbf{Vietoris--Rips
filtration} (\texttt{src/cloudforger/tda/filtration/rips.py}, using
\texttt{ripser}). For a point cloud $X$ with pairwise Euclidean distances
$d(x_i,x_j)$, the Rips complex at scale $r$ is
\[
  R_r(X) = \{\sigma \subseteq X : d(x_i,x_j) \le r \;\; \forall\, x_i,x_j\in\sigma\},
\]
and the filtration $\{R_r(X)\}_{r\ge0}$ induces persistent homology groups
$H_k(R_r(X))$ for $k=0,1$ (\texttt{maxdim=1}). Each connected topological
feature is summarized as a birth--death pair $(b,d)\in \mathbb R^2$; the
collection for dimension $k$ is the diagram $D_k = \{(b_i,d_i)\}$. This is a
single call, \texttt{ripser(cloud.points, maxdim=1)["dgms"]}, with no
subsampling — the full point set is used for the Rips complex, in contrast
to \texttt{raw\_pc}.

\subsection{Persistence images (\texttt{pi\_0}, \texttt{pi\_1})}

Each diagram $D_k$ is rasterized into a fixed $64\times64$ image
(\texttt{src/cloudforger/tda/vectorizer/persistence\_image.py}). Points are
first mapped from birth--death to birth--persistence coordinates,
$p_i = d_i - b_i$, then rendered as a weighted sum of isotropic Gaussians on
a regular grid $\{(u_a, v_b)\}_{a,b=1}^{64}$ over a fixed
\texttt{birth\_range} $\times$ \texttt{pers\_range} box:
\[
  I(u_a,v_b) \;=\; \sum_{i} w(p_i)\,
    \exp\!\left(-\frac{(u_a-b_i)^2+(v_b-p_i)^2}{2\sigma_{\mathrm{im}}^2}\right),
  \qquad w(p)=p \;\;(\texttt{linear\_weight}),
\]
with kernel bandwidth $\sigma_{\mathrm{im}}=0.05$ (config
\texttt{vectorization.sigma}), i.e.\ mild persistence-weighting that zeroes
out diagonal (zero-persistence) points and grows linearly with $p$. Each
homology dimension gets its own image and its own calibrated bounds
(\texttt{MultiChannelImager}): \texttt{birth\_range} and \texttt{pers\_range}
are \emph{not} inferred per-cloud but fixed once from the 99th percentile of
births/persistences pooled over the entire train/test diagram set
(\texttt{cloudforger.tda.calibration.calibrate}), so a pixel means the same
physical scale across every sample. The resulting array is flipped along the
persistence axis so persistence increases upward, and stored per-dimension
under \texttt{image\_tensors[dim]} with shape \texttt{(n\_clouds, 64, 64)}.
At load time \texttt{PersistenceImageDataset} adds the channel axis, so the
tensor entering the encoder is $(1,64,64)$.

\subsection{Betti curves (\texttt{betti\_0}, \texttt{betti\_1})}

The Betti curve counts, at each filtration scale on a fixed grid
$r_1<\dots<r_{128}\in[0,1]$ (\texttt{src/cloudforger/tda/features/betti\_curve.py}),
how many $k$-dimensional features are alive:
\[
  \beta_k(r_j) \;=\; \sum_i \mathbf 1\{\, b_i \le r_j < d_i \,\}.
\]
Pairs with an infinite death are dropped by default
(\texttt{drop\_infinite=True}); \texttt{normalize=False} means the curve is
left as a raw count rather than divided by the number of intervals. This
produces one length-128 vector per cloud per homology dimension,
concatenated when more than one dimension is requested
(\texttt{BettiCurveFeature.vector}); for a single-dimension experiment
(\texttt{betti\_0} or \texttt{betti\_1}) the encoder input is a length-128
vector.

\subsection{Pairwise / correlation features (\texttt{pairwise})}

This is a second-order spatial-statistics summary computed directly from the
point cloud, independent of persistent homology
(\texttt{src/cloudforger/stats/pair\_dist.py},
\texttt{PairDistanceCDF(CloudStatistic)}). Given the cloud's $N$ points, the
statistic draws $M=$\texttt{n\_samples} ($=5000$) random \emph{ordered pairs}
of distinct points $(i,j)$, $i\ne j$, uniformly with replacement, computes
their normalized Euclidean separation
\[
  r_m = \frac{\lVert x_i - x_j\rVert}{\mathrm{diam}(W)}, \qquad
  \mathrm{diam}(W) = \sqrt{\textstyle\sum_a (\text{high}_a-\text{low}_a)^2},
\]
(the region diagonal — a fixed property of the domain, not of the specific
cloud, so it does not leak per-cloud information into the normalization),
and evaluates the empirical CDF of $\{r_m\}$ on a fixed grid of 64 points in
$[0,1]$:
\[
  \hat F(r_j) = \frac{1}{M}\sum_{m=1}^M \mathbf 1\{r_m \le r_j\}.
\]
Implementation detail: values are sorted once and the grid is evaluated with
\texttt{np.searchsorted(values, grid, side="right")}, giving the same
$\hat F$ as the indicator sum but in $O(M\log M + G)$ instead of $O(MG)$. The
per-cloud RNG is seeded as \texttt{cloud.seed + 1{,}000{,}937}, so the
Monte-Carlo pair sample is reproducible per cloud but decorrelated from the
cloud-generation RNG stream. The 64-vector $\hat F$ is the feature
(\texttt{CorrelationFeatures.features["pair\_distance\_cdf"]}); with only one
statistic registered, \texttt{feature\_matrix} has shape
\texttt{(n\_clouds, 64)}.

\section{Feature encoding}

All encoders (\texttt{src/cloudforger/nn/encoders/}) subclass a common
\texttt{Encoder(nn.Module)} that fixes an output \texttt{embedding\_dim}
$=128$; the point is that every method, regardless of input shape, produces
a comparable 128-d embedding consumed by the same head (Sec.~\ref{sec:head}).

\subsection{Point cloud $\to$ embedding: PointNet-style encoder}

\texttt{PointNetEncoder} (\texttt{encoders/point\_cloud.py}) applies a
\emph{shared, per-point} MLP $\phi_\text{pt}:\mathbb R^{d}\to\mathbb
R^{256}$ (hidden widths $(64,128,256)$, ReLU after each linear layer) to
every one of the $n=512$ points independently — i.e.\ the same weights are
broadcast over the point dimension — and then applies a symmetric
(permutation-invariant) pooling operator, coordinate-wise max, to collapse
the point axis:
\[
  h_k = \phi_\text{pt}(x_k) \in \mathbb R^{256}, \qquad
  z = \max_{k=1,\dots,n} h_k \in \mathbb R^{256}, \qquad
  e = W_\text{proj}\, z + b_\text{proj} \in \mathbb R^{128}.
\]
Symmetric pooling (a simplified PointNet, without the input/feature
transform networks of the original architecture) guarantees the embedding
$e$ is invariant to the arbitrary order in which points are stored/sampled,
which matches the fact that a point cloud is a set, not a sequence.
Concretely, in code the batch tensor has shape \texttt{(B, n, d)}, the MLP
is applied via broadcasting over the last axis, and pooling is
\texttt{features.max(dim=1)}.

\subsection{Persistence image $\to$ embedding: small CNN}

\texttt{PIEncoder} (\texttt{encoders/persistence\_image.py}) treats each
$1\times 64\times 64$ image as a single-channel image and applies three
convolutional blocks, each halving spatial resolution:
\[
  \text{Conv}_{3\times3}(1{\to}32) \to \mathrm{ReLU} \to \mathrm{MaxPool}_2
  \to \text{Conv}_{3\times3}(32{\to}64) \to \mathrm{ReLU} \to \mathrm{MaxPool}_2
  \to \text{Conv}_{3\times3}(64{\to}128) \to \mathrm{ReLU},
\]
followed by global average pooling (\texttt{AdaptiveAvgPool2d(1)}) to a
$128$-vector independent of spatial size, then a linear layer
$W_\text{proj}\in\mathbb R^{128\times128}$ to the shared embedding space. Each
homology dimension ($H_0$, $H_1$) is trained as an entirely separate
single-channel model (\texttt{in\_channels=1}); the two are not fused in
this experiment (that fusion is what \texttt{models/multi\_modal.py} is for,
unused here).

\subsection{Betti curve / pairwise CDF $\to$ embedding: MLP (\texttt{StatsEncoder})}

Both fixed-length vector features (length-128 Betti curve, length-64
pair-distance CDF) share the same encoder class
(\texttt{encoders/stats.py}): a plain MLP with hidden widths $(128,128)$ and
ReLU activations, followed by a linear projection head:
\[
  h = \mathrm{ReLU}(W_2\,\mathrm{ReLU}(W_1 x + b_1) + b_2), \qquad
  e = W_\text{proj}\, h + b_\text{proj} \in \mathbb R^{128}.
\]
There is no architectural inductive bias here beyond a standard feed-forward
network — the vector is treated as an unordered/ordered feature bag, not as
a sequence or image.

\subsection{Optional covariates}

If \texttt{use\_covariates=True}, a per-cloud covariate matrix (shape
$(N,n_\text{cov})$, aligned with points) is reduced to its \emph{mean over
points} and concatenated \emph{after} encoding, directly onto the embedding
before the head: $e' = [\,e \,\Vert\, \bar c\,] \in \mathbb
R^{128+n_\text{cov}}$ (\texttt{CovariateHeadDataset},
\texttt{Experiment.run}). Covariates therefore bypass the encoder entirely
and are injected as raw statistics, not learned representations.

\section{Model head and loss function}
\label{sec:head}

\subsection{Label preprocessing}

Before training, the three raw parameters $\theta=(\lambda_p,\mu,\sigma)$
are transformed per-label
(\texttt{experiments/base.py}, \texttt{\_normalize\_labels\_by\_name}):
\begin{enumerate}
  \item an optional $\log$ transform (default: applied to \emph{every}
    label, since all process parameters here are positive scale-type
    quantities), $\tilde\theta_j=\log\theta_j$;
  \item per-parameter standardization using the mean/std computed over the
    \emph{full} train/test dataset (before the train/val/test split),
    $y_j = (\tilde\theta_j-m_j)/s_j$, with $s_j\!\to\!1$ if the empirical std
    is $0$.
\end{enumerate}
The network is trained to predict $y=(y_1,y_2,y_3)$, not $\theta$ directly;
inverting the transform at evaluation time is $\hat\theta_j =
\exp(m_j+s_j\hat y_j)$ (not shown above, but implied by storing
\texttt{label\_log\_mean}/\texttt{label\_log\_std}/\texttt{label\_transforms}
alongside the checkpoint for later reconstruction).

\subsection{Head architecture}

Every method funnels its 128-d (or $128{+}n_\text{cov}$-d) embedding through
the \emph{same} regression head class,
\texttt{ParameterEstimator(nn.Module)} (\texttt{nn/heads/paramest.py}):
\[
  \hat y \;=\; W_\text{out}\;
    \mathrm{Dropout}_{0.1}\!\big(\mathrm{ReLU}(W_h\, e' + b_h)\big) + b_\text{out}
  \;\in\; \mathbb R^{3},
\]
i.e.\ one hidden layer of width 64 (\texttt{hidden\_dims=(64,)}) with ReLU
and dropout $p=0.1$, then a linear output layer with $3$ units (one per
Thomas parameter; in general $n\_params=\texttt{labels.shape[1]}$, e.g.\ $2$
for the Matérn process). Encoder and head are composed by
\texttt{SingleModalModel} (\texttt{nn/models/single\_modal.py}):
\texttt{forward(x, covariates) = head(cat[encoder(x), covariates])}.

\subsection{Loss function and optimization}

The loss is plain mean squared error over the normalized targets,
\[
  \mathcal L(\hat y, y) \;=\; \frac{1}{B}\sum_{b=1}^{B}\lVert \hat y^{(b)} - y^{(b)}\rVert_2^2
  \qquad (\texttt{nn.MSELoss()}),
\]
i.e.\ every one of the (up to 3) regression targets is weighted equally in
the loss \emph{after} log/z-score normalization — since normalization
equalizes each label's variance to $1$, this keeps parameters on very
different natural scales (e.g.\ $\lambda_p\sim O(10\text{--}200)$ vs.\
$\sigma\sim O(0.005\text{--}0.1)$) from dominating the gradient.
Optimization (\texttt{nn/train.py}, \texttt{experiments/base.py
\_train\_and\_eval}):
\begin{itemize}
  \item optimizer: Adam, learning rate $10^{-3}$ (\texttt{cfg["lr"]});
  \item batch size 32, for 500 epochs (\texttt{cfg["batch\_size"]},
    \texttt{cfg["n\_epochs"]});
  \item data split: random $70/15/15$ train/val/test
    (\texttt{nn/splits.py: train\_val\_test\_split}, seeded by
    \texttt{cfg["seed"]});
  \item model selection: after every epoch, validation MSE is checked and
    the full \texttt{state\_dict} is checkpointed in memory whenever it
    improves; the \emph{best-val} weights (not the final-epoch weights) are
    reloaded before computing the reported test loss;
  \item an optional adversarial pass evaluates the frozen best model on a
    held-out set of parameter vectors disjoint from the train/test grid
    (\texttt{adversarial\_clouds.pkl} etc.), reusing the same label
    normalization statistics fit on the train/test set — this measures
    extrapolation to unseen $(\lambda_p,\mu,\sigma)$ combinations, not just
    unseen realizations of seen parameters.
\end{itemize}
Everything (best \texttt{state\_dict}, full model, loss curves, label
normalization constants, and — if run — adversarial loss) is serialized to
\texttt{results.pt}/\texttt{model.pt}/\texttt{results.json} under
\texttt{results/params/\{dim\}d/\{process\}/\{method\}/seed\_\{seed\}/}.

\end{document}
